


\section{Managerial Implications and Industry Perspectives}
\label{sec:implications}

Our findings extend beyond algorithmic validation to inform the design of decision support in next-generation Enterprise Resource Planning (ERP) systems. While conversational assistants can materially improve accessibility to data and routine workflows, our results suggest a practical limitation: when systems move from \emph{insight generation} to \emph{operational control}, reliability hinges less on linguistic fluency and more on explicit optimization, constraint enforcement, and verifiable execution. This motivates an architectural pattern in which the natural-language interface is separated from the decision engine, rather than treating a general-purpose LLM as an end-to-end controller.

\subsubsection*{Beyond the ``Chatbot'' Paradigm: Where Generalist AI Struggles.}
Recent market assistants---for example, \textit{Shopify Sidekick}, \textit{Square AI}, and \textit{Toast IQ}---have demonstrated strong capabilities in natural-language querying, summarization, and lightweight workflow automation (e.g., surfacing trends, drafting explanations, or supporting routine administrative actions). These features are valuable for \emph{descriptive} and \emph{diagnostic} analytics and can reduce the cognitive cost of interacting with complex business data.

However, extending the same paradigm to \emph{prescriptive} decision-making in stochastic operational settings introduces distinct failure modes. 
In our benchmark, the \emph{GPT-4o (Parameter Input)} baseline illustrates this point: even when provided with structured, high-fidelity state variables, a generalist LLM acting as the primary decision-maker did not reliably produce policies that matched the performance of optimization-based approaches. Under our modeled objective, this manifested as a substantial ``Heuristic Tax'' (exceeding $30\%$ in inventory cost in the evaluated setting). 
Moreover, recent field experiments demonstrate that autonomous agents tasked with revenue optimization can rapidly drift into ``hallucinated'' compliance—overriding price controls or ignoring profit constraints—when context limits are reached~\citep{wsj_vending_ai_2025}.
Importantly, this should not be interpreted as a universal statement about all deployed assistants; rather, it highlights a design risk: if an ``AI copilot'' is asked to generate \emph{executable} operational policies without an explicit solver, constraint checks, or systematic evaluation, it may produce recommendations that are plausible in narrative form but financially brittle in execution.

The managerial implication is therefore not to avoid LLM-based interfaces, but to scope them appropriately: use generalist models to \emph{elicit intent, clarify constraints, and explain trade-offs}, while delegating the computation of actions (and their validation) to specialized routines with auditable guarantees.






\subsubsection*{Lowering the Barrier to Operational Rigor.}
Advanced inventory policies (e.g., dynamic $(s, S)$ rules derived from dynamic programming or approximate stochastic control) have long been available in the academic and industrial toolkit. In practice, however, deploying these methods at scale has been disproportionately easier for firms with strong data infrastructure, analytics talent, and process discipline. Small and medium-sized enterprises (SMEs) often operate under constraints---limited analyst capacity, inconsistent master data, and tight operational bandwidth---that make it difficult to translate formal models into day-to-day decisions. As a result, many SMEs rely on stable heuristics (e.g., fixed reorder points or weekly manual adjustments), not because optimization is conceptually inaccessible, but because the \emph{implementation and maintenance costs} are nontrivial.

Our hybrid agentic framework targets this deployment gap. By acting as a semantic bridge, the system allows decision-makers to express business intent in natural language (service targets, ordering constraints, budget limits, and operational preferences) while executing the resulting decision problem with a solver-based backend. In this sense, the framework can \emph{productize} operational rigor: it reduces the expertise required to \emph{use} advanced policies, even though the policies themselves remain grounded in explicit formulations and validated computation.

The economic implications should be interpreted with appropriate caution and context. In our benchmark environment, the proposed approach achieved a 32.1\% reduction in the modeled inventory cost relative to heuristic baselines. If improvements of comparable direction translate to practice, they may materially affect cash flow and working-capital utilization, especially in sectors with thin margins and limited buffer capacity. That said, realized impact will vary with demand volatility, lead-time uncertainty, the baseline inventory-to-sales ratio, and execution frictions (data quality, supplier constraints, and organizational adherence). Managers should therefore treat performance gains as \emph{evaluable hypotheses} that require backtesting and staged rollouts rather than as guaranteed outcomes.

\subsubsection*{The Future of Intelligent ERP: A Three-Stage Design Pattern.}
We frame the evolution of management software as a design pattern with three stages:
\begin{itemize}
    \item \textbf{Type I (Descriptive):} Dashboards and reports that visualize data and leave decision logic entirely to human operators.
    \item \textbf{Type II (LLM-Augmented):} Systems that add conversational interfaces and agentic features to summarize data, draft recommendations, and trigger limited actions. Decision logic in these systems is often heuristic, partially opaque, or weakly constrained, which can be problematic for high-stakes operational control.
    \item \textbf{Type III (Orchestrating):} Systems in which the LLM is a \emph{front-end orchestrator} rather than the decision engine.
\end{itemize}

In a Type III system, the interface and the engine are intentionally decoupled. The LLM is used for (i) intent capture and constraint elicitation, (ii) translation into structured problem specifications, and (iii) explanation and what-if communication. Execution is delegated to specialized optimization and simulation routines, with guardrails such as constraint validation, audit logs, approval workflows for high-impact actions, and fallback policies when assumptions drift. This architecture does not eliminate uncertainty or model risk; instead, it makes the locus of decision-making explicit and testable.

We therefore view modular orchestration as a robust pathway toward ``prescriptive ERP'' in domains where decisions are frequent, costs are nonlinear, and uncertainty is material. The central managerial takeaway is to invest in systems that can \emph{communicate naturally} while remaining \emph{solver-validated under a specified formulation}, and to operationalize them with monitoring and governance rather than treating language fluency as a proxy for decision quality.



\section{Conclusion and Future Directions}
\label{sec:conclusion}

A central tension in management science lies between \emph{normative optimality}—the theoretical best prescribed by models—and \emph{descriptive reality}—the ambiguous, messy context in which practitioners operate. This paper proposes a Hybrid Agentic Framework to bridge this gap. By combining the semantic flexibility of Large Language Models with the rigorous precision of Operations Research solvers, we demonstrate that it is possible to democratize advanced inventory control without sacrificing mathematical fidelity.
Our empirical evaluation, conducted on a novel \emph{Human Imitator} testbed, confirms that this division of labor is superior to end-to-end LLM reasoning.


Our framework is designed as a modular foundation, inviting extension across several dimensions:

\begin{itemize}\item \textbf{Ecological Validation:} While our \emph{Human Imitator} provides a scalable proxy for bounded rationality, future work should validate the system against real-world human subjects. This involves refining the information extraction pipeline to handle the nuances, interruptions, and non-linear logic typical of actual manager-consultant dialogues.\item \textbf{Algorithmic and Agentic Expansion:} The current \emph{Optimization Agent} is agnostic to the underlying solver; future iterations could incorporate more advanced methodologies. Furthermore, we envision richer interaction patterns, such as a ``Critic Agent'' that detects inconsistencies between user goals and solver outputs, creating a verification loop before recommendations are presented.

\item \textbf{Scope Generalization:} Finally, this architecture can be extended beyond inventory management. By swapping the underlying solver, the framework could support other high-stakes decision domains—such as dynamic pricing, workforce scheduling, or logistics routing—where the synergy of natural language interfaces and rigorous optimization is critical.
\end{itemize}